%! AP Statistics — Unit 1, Lesson 1.5 — Student Packet Cover Sheet
\documentclass[10pt]{article}
\usepackage{apstats-article}
\usepackage{apstats-boxes}

%=============================================================================
\begin{document}

% ── Full-bleed navy banner ────────────────────────────────────────────────────
\begin{tikzpicture}[remember picture, overlay]
  \fill[navy] (current page.north west)
    rectangle ([yshift=-0.9in]current page.north east);
\end{tikzpicture}
\vspace{-0.8in}

\begin{center}
  {\color{white}\textbf{\LARGE AP Statistics}}\\[4pt]
  {\color{skymid}\large\textbf{Unit 1: Exploring One-Variable Data}}\\[6pt]
  {\color{white}\textbf{\large Lesson 1.5 \quad Representing a Quantitative Variable with Graphs}}
\end{center}

\vspace{0.22in}

\namedateperiod

\vspace{0.18in}

% ── LEARNING TARGETS ─────────────────────────────────────────────────────────
\begin{learningtargetbox}
\small
\begin{enumerate}[leftmargin=1.5em, itemsep=5pt]
  \item \ldots{} read and interpret a \textit{dotplot} for a quantitative dataset,
        identifying center, spread, shape, gaps, and clusters.
  \item \ldots{} complete and interpret a \textit{stemplot} (stem-and-leaf plot),
        including reading a back-to-back stemplot to compare two groups.
  \item \ldots{} read and interpret a \textit{histogram} from a frequency or
        relative frequency table, and explain why bars touch.
  \item \ldots{} identify which display (dotplot, stemplot, or histogram) is most
        appropriate for a given dataset and purpose.
\end{enumerate}
\end{learningtargetbox}

\vspace{0.14in}

% ── TABLE OF CONTENTS ────────────────────────────────────────────────────────
\begin{tocbox}
\small
\renewcommand{\arraystretch}{1.5}
\begin{tabularx}{\linewidth}{@{} c l X r @{}}
\rowcolor{sky}
\textbf{\#} & \textbf{Component} & \textbf{Description} & \textbf{Score} \\
\midrule
1 & Warm-Up        & Read a dotplot; compare to bar chart                & \blank{1.2cm} \\
2 & Guided Notes   & Dotplots, stemplots, histograms, trade-offs          & \blank{1.2cm} \\
3 & Group Activity & Analyze displays; complete a stemplot; interpret a histogram & \blank{1.2cm} \\
4 & Exit Ticket    & Read a stemplot; calculate median; describe shape    & \blank{1.2cm} \\
5 & Homework       & Dotplot reading, stemplot completion, AP-style       & \blank{1.2cm} \\
\midrule
  & \multicolumn{2}{r}{\textbf{Total}} & \blank{1.2cm} \\
\end{tabularx}
\end{tocbox}

\vspace{0.14in}

% ── KEEP IN MIND ─────────────────────────────────────────────────────────────
\begin{remindbox}
\small
Each display in this lesson answers a different question. Choose the right tool for the job.

\vspace{6pt}
\begin{tabularx}{\linewidth}{@{} >{\centering\arraybackslash\columncolor{goldbg}}p{0.06\linewidth}
                                    X @{}}
\textbf{\large!} &
\textbf{Histogram bars touch; bar chart bars do not.}\enspace
Bars touch in a histogram because adjacent bins share a boundary --- there is no
gap between consecutive intervals on a numerical scale. \\[6pt]
\textbf{\large!} &
\textbf{Every stemplot needs a key.}\enspace
Without a key, the reader cannot interpret the values.
Always write: e.g., \textbf{7 $|$ 4 = 74}. \\[6pt]
\textbf{\large!} &
\textbf{Use relative frequency, not count, when comparing histograms of different sizes.}\enspace
A taller bar in one histogram does not mean a higher proportion if the sample sizes differ. \\[6pt]
\textbf{\large!} &
\textbf{These graphs are the foundation of Lesson 1.6.}\enspace
The shape, center, spread, and outliers (SOCS) you will describe in 1.6 are all
visible in today's displays --- practice noticing them now. \\
\end{tabularx}
\end{remindbox}

\end{document}
